\Needspace{8\baselineskip}
\section{An implementation order with falsifiable acceptance gates}
\label{sec:roadmap}

\subsection{First gate: one complete weighted proof in Lean}

The weighted worker is the most economical first production slice. Its search
is elementary exact linear algebra, its certificate has three equations, and
its source fragment can initially be explicit rational matrices. Implement the
matrix closure theorem and a decoder that constructs dimension-indexed Lean
values. Then replay the stored subsequence-product certificate against an
explicit Lean definition of the two folds and a proved affine/tensor bridge.

The acceptance gate is not ``the matrices simplify''. It is a theorem over
\emph{all source words}, produced from the certificate, compiled in the pinned
environment, with the exact final theorem and axiom inventory retained. Change
one transition coefficient, omit a letter, and alter one closure entry; the
same replay path must reject each stale artifact. A distinguishing-word test
must produce a Lean-checked inequality on an actual source input.

Only after that gate should the source front end recognize a larger family of
folds automatically. Starting with explicit representation annotations isolates
semantic extraction from linear-algebra debugging and makes failures easier to
locate.

\subsection{Second gate: a non-conserved invariant from source code}

Port the invariant soundness theorem and the coefficient-expression replay
route. Use the nonlinear squaring curve as the minimum positive example, with
a natural-number iteration count and a rational state. The output theorem must
show the invariant for every iteration, not merely prove the one-step identity.

Require a second test with two possible transitions and a third with an equality
guard. Those prevent the implementation from accidentally specializing to
single-map conservation. Retain the degree-one finite-orbit miss as an unknown
control and the degree-two success as a separate acceptance. Do not make the
frontend's degree bound depend on seeing the target theorem's known proof.

For the CAS adapter, start with the included Python worker or the documented
Sage lifting interface. The final Lean proof should depend on the decoded
multipliers, not on which external search produced them. Changing the search
backend should require new performance tests, not a new trusted theorem.

\subsection{Third gate: the square sum with all endpoints}

Formalize the zero-extended binomial semantics or an equivalent piecewise
natural-index formulation. Prove the regularized summation theorem. Then use
the recorded square-sum certificate to prove the recurrence, prove the
candidate sequence $\binom{2n}{n}$ satisfies the same recurrence, and apply
uniqueness.

The negative gates should include the old rational expression evaluated at a
pole, a dropped endpoint contribution, a wrong range upper bound, and the
$(n-2)$-multiplied recurrence offered as an unrestricted uniqueness certificate.
The latter is especially valuable because its polynomial algebra is correct;
it tests semantic completeness of the acceptance obligations rather than
whether \code{ring} can detect a typo.

Next, replay the cube-sum recurrence and its two initial values. This acceptance
must remain labeled as a recurrence theorem unless a separate target sequence
has been proved to share the recurrence. A general-summation milestone should
add at least one nonzero-boundary example before expanding the recognized
summand grammar substantially.

\subsection{Controller integration and proof export}

Initially run each worker as a bounded child of the existing obligation
controller. Use the conditional-candidate path when source bridges remain.
Once a worker proves a useful consequence, assert that consequence as an
ordinary theorem and let the baseline tactic continue. An incremental
\grind{} extension is not necessary to validate the new mathematics.

Export should preserve three pieces: the source bridge theorem, the finite
certificate data, and the theorem applying the soundness principle to those
data. Ordinary recompilation should not need the external solver. A compact
local lemma is usually preferable to replaying search inside the final tactic
script.

For observability, attach worker-specific diagnostics rather than only a single
``failed'' counter. The ideal worker should report the degree, descent
dimensions, Gr\"obner budget status, and whether goal lifting failed. The
weighted worker should report ambient dimension, independent row count,
examined row count, and any distinguishing word. The telescoping worker should
report order, degree caps, coefficient-matrix size, nullity, and the first
unproved support or nonvanishing obligation. These are directly actionable
reasons to change a budget or repair a source bridge.

\subsection{Evaluation after the formalized slice exists}

Build a held-out Lean corpus grouped by the reason the worker should help:
nonconserved invariant families, independently implemented affine folds,
tensor-composed observations, constant-coefficient recurrences, and definite
hypergeometric sums with moving support. Include true unsupported goals and
false variants in each family. Record source theorem statements before tuning
the worker and exclude target theorems or equivalent lemmas from retrieval.

Compare the baseline tactic under the same imports, premises, and resource
limits with the baseline plus each worker individually and all workers
combined. Also compare against a simple portfolio of existing tactics, since
an advantage over a deliberately weak baseline would not establish the value
of the new algorithms. Record search time, elaboration time, kernel checking
time, proof size, solver startup, and replay-only cost separately.

The important ablation is semantic. Remove polynomial multipliers from the
ideal certificate language. Disable tensor composition from the weighted
frontend. Disable support regularization from the sum worker while retaining
its strict rejection of undefined obligations. These variants test the
incremental capabilities described in this report. Merely lowering a generic
search budget would not identify which mathematical mechanism caused success.

\subsection{Further directions that fit the same boundaries}

The most natural next invariant extension is a product of location-specific
spaces, followed by equality-aware treatment of piecewise polynomial updates.
For stronger semantic consequences, radical or real-radical certificates could
be added as a separate proof language; they should not be smuggled into ordinary
ideal membership.

For weighted behavior, add symbolic finite classifiers, several simultaneous
scalar observations, and proved reachable/observable minimization. For a vector
output, the certificate condition becomes $BB_{\mathrm{out}}=0$ with several
output columns, so the same closure search supports multiple contracts at
once. Polynomial outputs may be tensor-lifted selectively, but a dimension
budget is essential.

For sums, support inhomogeneous recurrences, exceptional initial parameter
ranges, and certificates expressed by a small operator basis rather than a
single hypergeometric ratio. A broader CAS backend should increase search
power while leaving the Lean summation obligations explicit.

These directions should follow, rather than postpone, the three complete Lean
acceptance gates. The present prototypes already provide enough exact data to
make those gates concrete.

\subsection{Assessment}

The proposed capability increase is specific. Forge would acquire a method for
finding invariant \emph{equations} that are not conserved quantities; a complete
finite-dimensional equivalence method for rational weighted behaviors; and a
method for deriving uniform recurrences for non-polynomial finite sums without
ignoring their support. Each returns small mathematical objects that ordinary
induction or finite summation can consume.

The most valuable next artifact is a compiled source-level theorem using one
of these certificates. The package makes that artifact substantially more
specified: the finite data, independent replay, mathematical soundness proof,
source fragment, relevant library interfaces, and rejection controls are all
present. What remains is the Lean implementation and its measurement, not an
unspecified request to make automation more powerful.
